% =============================================================================
% ALL PROMPT TEMPLATES
% =============================================================================
% Requires: \usepackage{listings} or use \begin{verbatim} / tcolorbox
% =============================================================================

\subsection{Prompt Templates}
\label{sec:prompts}

This section documents all prompt templates employed across the three extensions of the AskQE pipeline. For readability, placeholders are shown in curly braces.


% ----- PROMPT ABLATION -----

\subsubsection{Prompt Ablation (QA Stage)}
\label{sec:prompts_ablation}

The following three prompt variants modify the \textit{question answering} stage while keeping the question generation unchanged from the vanilla baseline.

\paragraph{P1: Few-Shot Medical.}
This variant provides two domain-specific in-context examples to guide the model toward concise, medically grounded answers.

\begin{quote}
\small\ttfamily
Answer medical questions based on the context.\\[4pt]
Example 1:\\
Context: "Patient diagnosed with diabetes mellitus type 2, prescribed metformin 500mg twice daily."\\
Question: "What medication was prescribed?"\\
Answer: "metformin 500mg"\\[4pt]
Example 2:\\
Context: "MRI revealed a 2cm lesion in the left frontal lobe."\\
Question: "What was the size of the lesion?"\\
Answer: "2cm"\\[4pt]
Now answer:\\
Context: \{sentence\}\\
Question: \{question\}
\end{quote}

\paragraph{P2: Chain-of-Thought.}
This variant instructs the model to reason step-by-step before producing a final answer, aiming to improve discrimination between severity levels.

\begin{quote}
\small\ttfamily
Answer the question about the medical text.\\[4pt]
Context: \{sentence\}\\
Question: \{question\}\\[4pt]
Think step by step:\\
1. Identify the key medical information in the context\\
2. Locate the specific answer to the question\\
3. Provide the exact answer from the text\\[4pt]
Format your response as:\\
REASONING: [your reasoning]\\
FINAL ANSWER: [concise answer using exact terms from context]
\end{quote}

\paragraph{P3: Concise Constraint.}
This variant enforces brevity to reduce verbosity and improve Exact Match alignment.

\begin{quote}
\small\ttfamily
Context: \{sentence\}\\
Question: \{question\}\\[4pt]
Answer in maximum 5 words using exact terms from the context:
\end{quote}


% ----- NER EXTENSION -----

\subsubsection{NER Extension (QG Stage)}
\label{sec:prompts_ner}

The NER extension modifies the \textit{question generation} stage by conditioning each question on a specific biomedical entity extracted via NER.

\paragraph{Entity-Aware Question Generation.}

\begin{quote}
\small\ttfamily
Task: Generate a specific question about the given entity based on the sentence.\\[4pt]
The entity is: "\{entity\}" (Type: \{entity\_type\})\\
The sentence is: "\{sentence\}"\\[4pt]
Generate exactly ONE clear, specific question that:\\
1. Can be answered using ONLY information from the sentence\\
2. Specifically asks about the entity mentioned\\
3. The answer should be the entity itself or information directly related to it\\[4pt]
Output only the question, nothing else.\\[4pt]
Question:
\end{quote}

\paragraph{Entity-Aware Question Answering.}

\begin{quote}
\small\ttfamily
Task: Answer the question based ONLY on the given sentence.\\[4pt]
Sentence: \{sentence\}\\
Question: \{question\}\\[4pt]
Instructions:\\
- Answer using ONLY information from the sentence\\
- If the answer is not in the sentence, respond with "[NOT FOUND]"\\
- Be concise and direct\\[4pt]
Answer:
\end{quote}


% ----- METRICS EXTENSION -----

\subsubsection{Metrics Extension (Evaluation Stage)}
\label{sec:prompts_metrics}

The metrics extension introduces two NLI-based classifiers at the \textit{evaluation} stage to assess semantic consistency between source and backtranslated answers.

\paragraph{LLM-as-Judge (Qwen2.5-3B-Instruct).}
The prompt constrains the model to produce a single label token, enabling first-token probability extraction (see Section~\ref{sec:nli_methodology}).

\begin{quote}
\small\ttfamily
You are a judge evaluating the relationship between two answers to the same question.\\[4pt]
Answer A (Source): \{answer\_src\}\\
Answer B (Backtranslation): \{answer\_bt\}\\[4pt]
Classify the relationship as one of:\\
- ENTAILMENT: Answer B supports or is consistent with Answer A\\
- NEUTRAL: Answer B is neither clearly supportive nor contradictory\\
- CONTRADICTION: Answer B contradicts or is inconsistent with Answer A\\[4pt]
Respond with ONLY the label: ENTAILMENT, NEUTRAL, or CONTRADICTION.\\
Label:
\end{quote}

\paragraph{NLI Classifier (BART-large-MNLI).}
The dedicated NLI model (\texttt{facebook/bart-large-mnli}) does not require a natural language prompt. It receives the source answer $a_s$ as premise and the backtranslated answer $a_{bt}$ as hypothesis via the standard \texttt{[SEP]}-separated input format:
\begin{equation}
    \texttt{input} = [a_s \;\texttt{</s>}\; a_{bt}]
\end{equation}
The model outputs logits over three classes (\textsc{Contradiction}, \textsc{Neutral}, \textsc{Entailment}), from which probabilities are obtained via softmax.
